% Rapport LaTeX (placeholders screenshots)
% Compile with: pdflatex rapport_final.tex (or latexmk)

\documentclass[12pt,a4paper]{article}
\usepackage[utf8]{inputenc}
\usepackage[T1]{fontenc}
\usepackage[french]{babel}
\usepackage{geometry}
\usepackage{graphicx}
\usepackage{hyperref}
\usepackage{caption}
\usepackage{float}
\geometry{margin=2.2cm}

% If screenshots are missing, show a placeholder instead of failing compilation.
\newcommand{\imgorplaceholder}[3]{%
  \begin{figure}[H]
    \centering
    \IfFileExists{#1}{\includegraphics[width=#2\textwidth]{#1}}{\fbox{\parbox[c][0.25\textheight][c]{#2\textwidth}{\centering Missing image\\\texttt{#1}}}}
    \caption{#3}
  \end{figure}
}

\title{Application de Gestion des PFE (ENSIT) \\ Rapport}
\author{\textit{Nom Prénom}}
\date{\today}

\begin{document}
\maketitle
\tableofcontents
\newpage

\section{Introduction}
Ce mini-projet consiste à développer une application Java (Swing + JDBC) pour gérer les Projets de Fin d'Étude (PFE) : étudiants, encadreurs, PFE et soutenances.

\section{Technologies}
\begin{itemize}
  \item Java (Swing)
  \item JDBC
  \item MySQL
  \item IDE : NetBeans
\end{itemize}

\section{MCD}
\subsection{Entités}
\begin{itemize}
  \item Etudiant(idEtudiant, nom, prenom, email, classe)
  \item Encadreur(idEncadreur, nom, prenom, grade, email)
  \item PFE(idPFE, titre, description, etat, dateSoutenance)
  \item Soutenance(idSoutenance, date, salle, note)
\end{itemize}

\subsection{Cardinalités}
\begin{itemize}
  \item Etudiant (0..1) --- (1..1) PFE
  \item PFE (1..n) --- (0..n) Encadreur
  \item PFE (0..1) --- (1..1) Soutenance
\end{itemize}

\section{MLD}
\begin{itemize}
  \item etudiant(id\_etudiant PK, nom, prenom, email UNIQUE, classe)
  \item encadreur(id\_encadreur PK, nom, prenom, grade, email UNIQUE)
  \item pfe(id\_pfe PK, titre, description, etat, date\_soutenance NULL, id\_etudiant FK UNIQUE)
  \item pfe\_encadreur(id\_pfe FK, id\_encadreur FK, PK(id\_pfe,id\_encadreur))
  \item soutenance(id\_soutenance PK, date, salle, note NULL, id\_pfe FK UNIQUE)
\end{itemize}

\section{Implémentation (résumé)}
\subsection{Structure simple}
\begin{itemize}
  \item \texttt{model} : classes métier (POJO)
  \item \texttt{dao} : CRUD JDBC (PreparedStatement)
  \item \texttt{ui} : fenêtres Swing (menu + modules)
\end{itemize}

\subsection{Contraintes et validations}
\begin{itemize}
  \item Email : format valide (étudiants/encadreurs)
  \item PFE : 1 étudiant \rightarrow max 1 PFE
  \item PFE : au moins 1 encadreur
  \item Soutenance : salle = A ou B uniquement
  \item Note : optionnelle, si saisie entre 0 et 20
  \item 1 PFE \rightarrow max 1 soutenance
\end{itemize}

\subsection{Sécurité minimale (recommandée)}
Créer un utilisateur MySQL dédié (au lieu de root) avec droits limités à \texttt{gestion\_pfe}.

\section{Captures d'écran (placeholders)}
Place les images dans le dossier \texttt{screenshots/} (au même niveau que ce fichier .tex).

Liste et ordre recommandés : \texttt{screenshots/SCREENSHOTS\_LIST.md}.

\imgorplaceholder{screenshots/01_menu.png}{0.9}{Menu principal}

\imgorplaceholder{screenshots/02_etudiants_liste.png}{0.9}{Gestion des étudiants}

\imgorplaceholder{screenshots/03_etudiants_validation_email.png}{0.9}{Validation email (étudiants)}

\imgorplaceholder{screenshots/04_encadreurs_liste.png}{0.9}{Gestion des encadreurs}

\imgorplaceholder{screenshots/05_encadreurs_validation_email.png}{0.9}{Validation email (encadreurs)}

\imgorplaceholder{screenshots/06_pfe_liste.png}{0.9}{Gestion des PFE}

\imgorplaceholder{screenshots/07_pfe_multi_encadreurs.png}{0.9}{Sélection multiple des encadreurs (PFE)}

\imgorplaceholder{screenshots/08_pfe_etudiant_unique.png}{0.9}{Contrainte : un seul PFE par étudiant}

\imgorplaceholder{screenshots/09_soutenances_liste.png}{0.9}{Gestion des soutenances}

\imgorplaceholder{screenshots/10_soutenances_salle_AB.png}{0.9}{Validation salle (A/B uniquement)}

\imgorplaceholder{screenshots/11_soutenances_note.png}{0.9}{Validation note (0..20)}

\section{Jeux de tests}
Le script \texttt{seed\_data.sql} permet d'insérer des exemples (étudiants/encadreurs/PFE/soutenance).

\section{Conclusion}
L'application fournit une gestion complète (CRUD) des étudiants, encadreurs, PFE et soutenances avec validations basiques.

\end{document}

